% Generado por lima/tablas_lima.py a partir de lima/salidas/ y lima/alta/salidas/. No editar a mano.
% Caso Lima, prueba de robustez del reparto: retraso total (s/veh) de cada control con el mismo total
% (4271 veh/h) repartido de otra forma entre las entradas, sin reajustar ni reentrenar, y cambio del modelo
% frente a cada control (%). 100 semillas (24001-24100); la primera fila es la tanda 1 del caso.
% fuente: lima/robustez/<variante>/salidas/eval_t1_*.csv; IC95 y p: reducto: actual -26.46 [-28.22, -24.78] p=3.9e-18; fijo -36.07 [-39.21, -32.68] p=3.9e-18; actuadoaj +4.04 [+3.43, +4.70] p=3e-17 | transversal: actual -16.32 [-17.22, -15.43] p=3.9e-18; fijo -29.94 [-31.95, -27.94] p=3.9e-18; actuadoaj +3.08 [+2.67, +3.51] p=4.5e-17
\begin{tabular}{>{\raggedright\arraybackslash}p{2.6cm}>{\raggedleft\arraybackslash}p{1.45cm}>{\raggedleft\arraybackslash}p{1.45cm}>{\raggedleft\arraybackslash}p{1.45cm}>{\raggedleft\arraybackslash}p{1.45cm}>{\raggedleft\arraybackslash}p{1.45cm}>{\raggedleft\arraybackslash}p{1.45cm}>{\raggedleft\arraybackslash}p{1.45cm}}
\toprule
Reparto de la demanda & Plan actual & Fijo ajustado & Actuado ajustado & Modelo completo & Frente al plan actual (\%) & Frente al fijo (\%) & Frente al actuado (\%) \\
\midrule
El supuesto (tanda 1) & 59.1 & 57.8 & 53.4 & 55.3 & $-6.5$ & $-4.4$ & $+3.4$ \\
Más carga en Reducto & 76.5 & 88.0 & 54.1 & 56.3 & $-26.5$ & $-36.1$ & $+4.0$ \\
Más carga en 28 de Julio y Balboa & 64.1 & 76.6 & 52.1 & 53.7 & $-16.3$ & $-29.9$ & $+3.1$ \\
\bottomrule
\end{tabular}
